\chapter*{Conclusion: What FRFP Learns from Its Case Studies}
\addcontentsline{toc}{chapter}{Conclusion: What FRFP Learns from Its Case Studies}

Volume~II has treated a diverse set of systems as FRFP objects:
trading desks and oncology clinics; hiring engines and risk scores; vending machines and robotaxis; protein-structure predictors and legal chatbots.
Despite the diversity, a small number of structural themes recur with striking regularity.
This conclusion distils those themes into a set of lessons for designers, institutions, and theorists.

\section*{Lesson 1: Explicit Success Does Not Guarantee Tacit Safety}

Many of our failure case studies looked superficially impressive in explicit metrics:

\begin{itemize}
  \item Watson for Oncology passed internal benchmarks and demos;
  \item Zillow's pricing models performed well on backtests;
  \item Amazon's recruiting engine optimised predictive accuracy on historical data;
  \item COMPAS satisfied certain calibration metrics;
  \item the CIO's VaR model produced numbers within nominal limits.
\end{itemize}

Yet in each case, the explicit semantics were misaligned with tacit notions of correctness:

\begin{itemize}
  \item Watson's training on synthetic cases and idiosyncratic expert judgements did not capture the real distribution of oncology cases;
  \item Zillow's models ignored selection effects and regime shifts that local experts understood;
  \item COMPAS could not simultaneously meet all fairness desiderata in the presence of base-rate differences;
  \item VaR metrics failed to encode key tail risks and liquidity constraints.
\end{itemize}

By contrast, AlphaFold and DR screening AI succeed precisely where explicit semantics are strong and well-matched to tacit needs.
The mapping from sequence to structure, or from retinal image to a clinically defined threshold, is compact enough to be captured by explicit models, and ground truth is abundant.

FRFP’s first conclusion is therefore:

\begin{quote}
  \emph{Explicit performance metrics are not proxies for tacit safety or correctness.
  They are informative only when the explicit task is well-specified, the data and model semantics match the tacit function we care about, and the limits of that match are understood.}
\end{quote}

\section*{Lesson 2: Boundaries are Designed, Not Assumed}

A central claim of Volume~I is that correctness lives at boundaries:
at the points where explicit artefacts are endorsed, rejected, or transformed by agents in tacit space.
Volume~II shows how often these boundaries are treated as rhetorical rather than structural.

In many failures:

\begin{itemize}
  \item systems are labelled ``decision support'' while functioning as de facto decision-makers (COMPAS, recruiting engines, recommenders);
  \item lawyers or clinicians treat AI outputs as if they were vetted results (Watson, \emph{Mata}) rather than hypotheses;
  \item risk metrics and model outputs replace human assessment at the point of commitment (Zillow offers, London Whale limits).
\end{itemize}

In successful or more robust systems:

\begin{itemize}
  \item DR AI explicitly reverts to human care when images are ungradable;
  \item AlphaFold outputs are clearly presented as predictions with confidence, not as experimental facts;
  \item Waymo draws a sharp operational boundary via its ODD, with human and institutional sign-off governing any expansion.
\end{itemize}

Our second conclusion:

\begin{quote}
  \emph{Boundaries between explicit systems and tacit endorsement must be deliberately engineered and institutionally enforced.
  Naming a system ``supportive'' or ``advisory'' does not make it so; interfaces, logging, and governance must make boundary events concrete and auditable.}
\end{quote}

The first attempt at long-horizon collaboration failed partly because it took boundaries for granted:
it assumed that humans would ``stay in the loop'' without describing the mechanisms that would force this to remain true under pressure.

\section*{Lesson 3: Safe Horizons are Structural, Not Optional}

The notion of a safe horizon---a bound on how far a run may proceed without re-grounding in tacit judgement---is perhaps the most distinctive FRFP contribution.
Volume~II confirms that, across domains, trouble arises where runs are effectively unbounded:

\begin{itemize}
  \item trading strategies accumulate risk over months without strong episodic review (London Whale, Zillow);
  \item long-horizon judicial use of risk scores gradually reshapes practice (COMPAS);
  \item repeated use of LLMs and legal-research surrogates erodes verification norms (\emph{Mata} and its successors);
  \item recommender systems adjust millions of people's content diets continuously without clear reset mechanisms.
\end{itemize}

Successful systems encode safe horizons in their very design:

\begin{itemize}
  \item DR programmes enforce periodic re-screening and post-market validation;
  \item Waymo restricts autonomy to defined ODDs and phases in expansions;
  \item scientific practice around AlphaFold, by tradition, requires experimental replication and peer review.
\end{itemize}

Thus:

\begin{quote}
  \emph{Any long-horizon deployment of AI---clinical, financial, legal, or epistemic---needs explicit safe horizons:
  limits in time, scale, and scope, together with procedures for systematic re-grounding.
  Without them, even well-behaved systems will drift into regimes where their explicit semantics no longer match tacit reality.}
\end{quote}

The initial, informal human--AI collaboration framework lacked such a notion.
It described workflows but not their depth limits.
FRFP arose when it became clear that we needed not just good workflows but principled bounds on their iteration.

\section*{Lesson 4: Governance Itself Cannot Be Automated}

Maybe the deepest pattern in Volume~II is the failure of explicit-only governance.

In case after case:

\begin{itemize}
  \item boards and executives rely on risk dashboards, backtests, and demos to justify scaling (Watson Health, Zillow, CIO);
  \item courts and regulators accept black-box scores under the rubric of ``data-driven'' decision-making (COMPAS, early risk tools);
  \item firms treat AI tools as internal experiments without firm-level policies (Amazon recruiting, early generative-AI use in law);
  \item platforms rely on engagement and safety metrics to manage recommender systems without fully grappling with their societal externalities.
\end{itemize}

Where governance is more successful:

\begin{itemize}
  \item regulators explicitly define indications, performance thresholds, and post-market surveillance for DR AI;
  \item scientific communities, through peer review and replication, institutionalise scepticism toward explicit models, even AlphaFold;
  \item autonomous-vehicle regulators and cities negotiate permits, ODD constraints, and incident-reporting frameworks with operators like Waymo.
\end{itemize}

This supports FRFP’s non-automation claims at the institutional level:

\begin{quote}
  \emph{There is no purely explicit procedure that, given model outputs and metrics, can reliably determine whether a socio-technical system is acceptable.
  Governance and oversight are themselves tacit functions carried out by humans and institutions, and must be modelled as such.
  Good governance uses explicit tools and metrics, but does not reduce itself to them.}
\end{quote}

The first collaboration framework implicitly treated governance as a design parameter: write down the right rules and processes, and the system will behave.
FRFP, informed by the case studies here, insists that governance is dynamic, contested, and structurally non-automatable, and therefore must be part of the formal semantics.

\section*{Lesson 5: FRFP is Both Descriptive and Prescriptive}

A final lesson concerns the role of FRFP itself.

On the descriptive side, Volume~II demonstrates that a single vocabulary---explicit/tacit spaces, runs, boundaries, safe horizons, institutional operators---is capable of reorganising a wide range of case studies into a small number of recurring structural patterns.
This is evidence that FRFP is not merely an arbitrary formalism but is tracking real invariants in human--AI systems.

On the prescriptive side, the case studies suggest concrete design moves:

\begin{itemize}
  \item \textbf{Narrow and stabilise explicit tasks} wherever possible, especially in high-stakes domains, and be honest about where semantics become diffuse (e.g.\ general-purpose LLMs).
  \item \textbf{Make uncertainty explicit} in interfaces, including ``I don't know'' and ``ungradable'' states that revert to human judgement.
  \item \textbf{Elevate boundaries to first-class objects}:
        log them, design UI around them, and give institutional structures authority at those points.
  \item \textbf{Encode safe horizons} in policies:
        caps, phases, revalidation triggers, and reset protocols, rather than relying on ad hoc audits after failures.
  \item \textbf{Build governance as an explicit layer} with its own semantics and operators:
        cross-functional committees, regulatory regimes, professional norms, and public deliberation.
\end{itemize}

FRFP does not guarantee good outcomes, any more than having a theory of mechanics guarantees a bridge will not fall.
But just as mechanical engineering makes certain failure modes legible (buckling, resonance, fatigue), FRFP makes certain socio-technical failure modes legible and therefore, in principle, avoidable.

\section*{Looking Forward: From Volume II Back to Practice}

The story of this monograph is itself a small FRFP case study.
An initial, mostly narrative attempt to design long-horizon human--AI collaboration proved inadequate.
That failure, when reflected on, led to a more formal architecture (Volume~I).
The architecture, in turn, has been tested against real systems (Volume~II) and refined by their successes and failures.

The next step is to close the loop:

\begin{itemize}
  \item to use FRFP prospectively in the design of new systems and institutions;
  \item to co-develop mathematical results, engineering practices, and governance patterns;
  \item to treat FRFP not as a static theory of ``AI safety'' but as an evolving language shared across technical, social-science, and institutional actors.
\end{itemize}

The case studies in this volume show that powerful AI can be embedded in human institutions in ways that are disastrous, benign, or genuinely transformative.
They also suggest that the difference is not mysterious:
it lies in the structure of the explicit pipelines, the design of boundaries and safe horizons, and the seriousness with which we treat institutional non-automation.

If Volume~I argued that such structure is \emph{necessary}, Volume~II suggests that it is \emph{already} present, implicitly, in systems that work reasonably well, and conspicuously absent in those that fail.
The task ahead is to make that structure explicit---in our tools, our organisations, and our laws---before the next generation of long-horizon AI systems arrives by default rather than by design.
